\TC{tea-club-01.mp3}{\TCtitle{\#1}{Fiola Soft}}{2014-04-10}{\Luke, \Fiola}{S Filipem Kraucherem s Fiola Softu si povídáme hlavně o čerstvě oznámené plošinovce Blackhole.}{0:29:26}{14964584}{}
\TC{tea-club-02.mp3}{\TCtitle{\#2}{Hyperbolic Magnetism}}{2014-05-08}{\Luke, \JIlavsky, \VHrincar}{Kdyby existovala statistika o počtu vyprodukovaných her za rok, programátorské duo Hyperbolic Magnetism by se s přehledem umístilo na prvních místech. Vladimír „Loki“ Hrinčár a Honza „Split“ Ilavský programují jak o život a pravidelně se hlásí na vývojářské maratony Global Game Jam a Ludum Dare. S jejich tituly se nejčastěji můžete setkat na mobilech, ale i na PC, které ve své hyperproduktivitě nezanedbávají. Příklady z poslední doby mohou být třeba Heterochroma z Global Game Jam nebo Honzova samostatná Temple of Ra z Ludum Dare.}{0:34:15}{19255532}{}
\TC{tea-club-03.mp3}{\TCtitle{\#3}{}}{2014-05-24}{\Luke, \JPavlasek}{O strategické tahovce Cubesis, ve které ovládáte ekosystém planety, jste možná neslyšeli. Nic si z toho nedělejte, my také ne. Proto jsme do dalšího Tea Clubu pozvali jednoho ze spoluautorů, aby nám o hře, jejíž vývoj trval sedm let, něco pověděl.}{0:29:23}{17355276}{}
\TC{tea-club-04.mp3}{\TCtitle{\#4}{Cinemax}}{2014-06-03}{\Luke, \LMacura, \MLinda}{Lukáš Macura a Martin Linda vyprávějí o vzniku konceptu a vymýšlení hádanek a soubojů pro krokovací dungeon The Keep na Nintendo 3DS a PC. Kromě toho jsme povídají o situaci na 3DS trhu a hledání toho správného publika. Probrali jsme také, jak se testují hry nejen v Cinemaxu, ale i v široširém světě.}{0:22:00}{11567484}{}
\TC{tea-club-05.mp3}{\TCtitle{\#5}{Future Factory}}{2014-06-07}{\Luke, \VGersl}{Po zkušenostech s prací v herním vývoji založil Vláďa Geršl studio Fun 2 Robots a připravuje hru, která se snad bude líbit všem příznivcům taktických akcí a literatury. Že to nejde dohromady? Future Factory obsahuje svižnou akci a příběh inspirovaný Čapkovým R.U.R. }{0:24:43}{12749724}{}
\TC{tea-club-06.mp3}{\TCtitle{\#6}{Novus Inceptio}}{2014-06-30}{\Luke, \MMastny}{MMORPG Novus Inceptio s procedurálně generovanou krajinou nenabídne tisíce naskriptovaných úkolů a předem připravená města a vesnice osázená vesničany. Herní obsah budou tvořit samotní hráči svojí aktivitou, která začíná přežívání v divočině a končí třeba stavbou města, volením starosty a pronajímáním pozemku.}{0:26:43}{13034020}{}
\TC{tea-club-07.mp3}{\TCtitle{\#7}{}}{2014-07-03}{\Luke, \MJezek}{Ve studiu 2K prý měli plyšového ježka zavřeného v lahvi. Když se jí Matouš Ježek pokusil otevřít, rychle ho zastavili s tím, že ježek se nikdy nesmí dostat ven. Jak to dopadlo s plyšákem nevíme, ale cestu Matouše ze studia 2K můžeme dobře sledovat. Po odchodu založil v Brně studio Trickster Arts a loni vydal pozitivně přijatou mobilní hru Hero of Many, která vyšla i pro OUYA a Amazon Fire a v současnosti se připravuje i pro Windows.}{0:24:43}{11794000}{}
\TC{tea-club-08.mp3}{\TCtitle{\#8}{Factorio}}{2014-08-03}{\Luke, \MKovarik, \TKozelek}{Michal Kovařík a Tomáš Kozelek nám v půl hodinovém rozhovoru popsali celou cestu jejich projektu Factorio – od prvních krůčků, zoufalství a naivity až k funkčnímu „garážovému“ týmu, který jde cílevědomě za svou vizí. Kluci mluví i o vlídném přijetí komunitou anebo o běhání jako očistném procesu v dobách, kdy jde vše do háje.}{0:30:57}{15993392}{}
\TC{tea-club-09.mp3}{\TCtitle{\#9}{Překlady her do češtiny}}{2014-10-25}{\Luke,\,\dots}{V Tea Clubu jsme se tentokrát zaměřili na slasti a strasti amatérského překládání her, konkrétně Divinity: Original Sin. Michal a Vítek, kteří šéfují týmu překladatelů a vedou komunitu na sociálních sítí, mluví o náročnosti překladu ukecaného RPGčka.}{0:16:48}{8066480}{}
\TC{tea-club-10.mp3}{\TCtitle{\#10}{SIMT MHD}}{2014-12-09}{\Luke, \TFaina}{Každý den hráči najezdí v autobusech a trolejbusech v SIMT téměř dva tisíce kilometrů, nabírají a vykládají cestující a snaží se dodržovat jízdní řád. Podle Tomáše Fainy, který nám o hře povídal v Tea Clubu, jim to docela jde a jsou kultivovaní. No bus trolling in da štrýt! Uznejte sami – v kolika hrách vám cestující nenastoupí, když máte zasviněný autobus? Takže pěkně hadru a Savo a drhněte!}{0:19:20}{9485636}{}
\TC{tea-club-11.mp3}{\TCtitle{\#11}{Wayward Terran Frontier}}{2014-12-29}{\Luke, \JOrszulika}{Wayward Terran Frontier je projektem Američana George Hultgrena a Ostravaka Honzy Orszulika, který se podílí na vzhledu hry a herních pravidlech. Na první pohled jde sice o jednoduchou 2D hru ve vesmíru, ale jakmile se podíváte zblízka, nestačíte se divit: rozlehlý vesmír, design interiéru vlastní lodi a základny, abordace nepřátel, přestřelky v koridorech, destrukce jednotlivých částí nepřátelského vybavení a spousta dalšího.}{0:16:26}{7809584}{}
\TC{tea-club-12.mp3}{\TCtitle{\#12}{Medieval Engineers}}{2015-01-25}{\Luke, \MRosa}{Marek Rosa ze studia Keen Software House (Miner Wars 2081, Space Engineers) vypráví o složitostech při vytváření středověké inženýrské simulace a také zmiňuje spoustu možností, které hru možná čekají.}{0:28:17}{13337132}{}
\TC{tea-club-13.mp3}{\TCtitle{\#13}{The Great Wobo Escape}}{2015-04-11}{\Luke, \RKajan, \MWilczak}{Pavel si tentokrát pod povídal s šéfem studia Gamifi.cc Rudou Kajanem a s grafikem Martinem Wilczakem, jak jejich hru The Great Wobo Escape inspirovaly klasiky jako Abe's Odyssee, Another World nebo film Wall-E. Dozvíte se ale samozřejmě mnohem víc.}{0:18:35}{8060756}{}
\TC{tea-club-14.mp3}{\TCtitle{\#14}{Capricorn 7}}{2015-04-26}{\Luke, \DKos}{Možná by si Future Cop: LAPD zasloužila zařadit mezi retro harampádí. Vyšla před necelými dvaceti lety na PlayStation a dnes vypadá stejně jako všechny rané 3D hry: staře a nekoukatelně. Kdo si na ní pamatuje, ten ale nehodnotí grafiku, nýbrž vynikající městskou akci s mechwarriorem, kterou Electronic Arts dokázala přinést jak v kampani, tak v multiplayeru. David Kos šel ještě dál. Podle vzoru Future Cop vystavěl v Unity vlastní Capricorn 7 a doufá, že uspěje na IndieGoGo.}{0:23:53}{10284908}{}
\TC{tea-club-15.mp3}{\TCtitle{\#15}{Rememoried}}{2015-07-12}{\Luke, \VKoudelka}{Jestli jste před viděli trailer na českou hru Rememoired, možná jste se podivili. Co to má být? To je hra? A o čem? Prý se v ní bude pracovat se sny, tvrdila tisková zpráva. Z toho jsme nikdo nebyli moudřejší. Proto v novém Tea Clubu přijal pohovku pro hosta Vláďa Koudelka, který jako samojediný autor za Rememoired stojí.}{0:22:35}{9876452}{}
\TC{tea-club-16.mp3}{\TCtitle{\#16}{Jindřich Skeldal}}{2015-11-17}{\Luke, \Skeldal}{Brány Skeldalu: 7 mágů sice čerpá jednak z klasických dungeonů jako Wizardry, Legend of Grimrock nebo Eye of the Beholder, ale mytologie hry je ovlivněna další Jindřichovou tvorbou. Kromě her jsou to i knihy, z nichž více jak tisíci stránkový román Brány Skeldalu čeká na vydání, nebo legendární westerny jako Sedm statečných. Z naší půl hodinky se toho zkrátka dozvíte hodně.}{0:34:14}{14997884}{}
\TC{tea-club-17.mp3}{\TCtitle{\#17}{Logistics}}{2016-01-23}{\Luke, \JZeleny}{Hrajete Open Transport Tycoon Deluxe a je vám líto, že nevychází moderní hry, které by na tradici dopravních tycoonů navazovaly? Neplačte! Logistics od českého programátora Honzy Zeleného představuje lákavou alternativu.}{0:34:56}{18110478}{}
\TC{tea-club-18.mp3}{\TCtitle{\#18}{Planet Nomads}}{2016-02-28}{\Luke, \DMaslovsky, \PRuzicka}{Kombinace stavění domů, ale i vozidel, prozkoumávání procedurálně generované planety a přežívání mezi agresivními zvířaty přinesla ostravskému studiu Craneballs téměř sto tisíc liber na Kickstarteru. O obsahu, a také o komplikacích vývoje, přišli do našeho studia popovídat Dan Maslovský a Petr Růžička, pro které je Planet Nomads jednak úspěšným projektem, a pak premiérou na PC.}{0:31:03}{13062142}{}
\TC{tea-club-19.mp3}{\TCtitle{\#19}{The Fall of the Fly}}{2016-04-24}{\Luke, \JValenta}{I když se snažím poctivě a metodicky kopírovat, tak mi z toho vzniká něco bláznivého, říká v jednu chvíli Honza Valenta, autor projektu Fall of the Fly a zároveň host dalšího Tea Clubu. Jeho hra je na první pohled skutečně... jiná. Někdo by možná řekl šílená a nesmyslná, ale i to je záměr. Přeci si nemyslíte, že uprostřed Dutozemě by měl být normální svět obydlený běžnými smrtelníky?}{0:20:00}{9596480}{}
\TC{tea-club-20.mp3}{\TCtitle{\#20}{Vaporum}}{2016-05-22}{\Luke, \TRepta, \MZajacik}{Je tomu už více než rok, co jsme psali o slovenském krokovacím dungeonu Vaporum. Proniknete v něm do vysoké věže s vnitřnostmi z parních strojů a arteriemi z párou natlakovaného potrubí. K čemu slouží? Odpověď znají mechanické příšery postávající na každém rohu, ale ty vám jí neřeknou. Jste pro ně vetřelec, lidský ztroskotanec, který se shodou okolností stává hrdinou. Kromě příběhu klade Vaporum důraz na neobvyklé vylepšování hrdiny a souboje, které se mají vymykat (nejen) moderním krokovacím klasikám jako Legend of Grimrock. Takže: Držte si gombíky, jdeme krokovat!}{0:39:37}{18375116}{}
\TC{tea-club-21.mp3}{\TCtitle{\#21}{The Cycle}}{2016-06-01}{\Luke, \MBudik}{Přibližně čtyři roky pracuje Marek Budík na střílečce The Cycle, která vypadá jako by vypadla z geometrova šíleného snu. Procedurálně generované úrovně jsou tu osazené nepřáteli, kteří čekají na příchod hrdiny s velkým mečem, aby ho zabili – nebo byli zabiti. Rychlé reakce jsou pro přežití výhodou, hrát se však dá i pomaleji a s rozmyslem. Na The Cycle kromě svižné akce potěší i grafická stylizace, která odkazuje k minimalismu a geometrickému vidění světa. Marek v herním návrhu odráží i to, co se naučil během studií průmyslového designu, a nebojí se o tom povídat.}{0:20:24}{9388904}{}
\TC{tea-club-22.mp3}{\TCtitle{\#22}{Clayburger o Hopply}}{2016-06-11}{\Luke, \PKnut, \MPavelka}{Pamětníci budou hned doma: Vždyť jde o Perestrojku, fenomén učeben informatiky v první polovině devadesátých let! Peter a Matěj inspiraci nezastírají. Jde jim především o to udělat perfektní skákací rychlík, který perfektně funguje a má výraznou výtvarnou stránku. Jak se to dělá? I to uvidíte na vlastní oči během našeho nového Tea Clubu o hře Hopply.}{0:21:45}{9913208}{}
\TC{}{\TCtitle{\#23}{Battle for Nethervein}}{2016-08-19}{\Luke, \JVanecek}{Honza Vaněček a jeho studio From the Void rozhodně nejsou začátečníci. Honza je české deskovkové komunitě dobře znám svojí karetní hrou Shuffle Heroes a předchozí McJohny's. V našem rozhovoru se kromě Battle for Nethervein rozpovídal o své minulosti a o deskovkové komunitě v Česku, i o náročnosti při prosazování českých deskovek v zahraničí či o komplikacících distribuce anebo o pořádání deskovkových game jamů.}{0:20:24}{}{}
\TC{tea-club-24.mp3}{\TCtitle{\#24}{Morning Men}}{2016-10-03}{\Luke, \LNizky}{Když se podíváte na obrázky ze slovenské adventury Morning Men, hodí vám do očí písek z pouštní Duny a na solar udeří direktem od M\oe{}bia. A není to náhoda. K jeho práci se výtvarníci ze studia Pixel Federation hrdě hlásí -- a ne pouze k němu\,\dots}{0:19:18}{9888260}{}
\TC{tea-club-24.mp3}{\TCtitle{\#24}{Morning Men}}{2016-10-03}{\Luke, \LNizky}{Když se podíváte na obrázky ze slovenské adventury Morning Men, hodí vám do očí písek z pouštní Duny a na solar udeří direktem od M\oe{}bia. A není to náhoda. K jeho práci se výtvarníci ze studia Pixel Federation hrdě hlásí -- a ne pouze k němu\,\dots}{0:19:18}{9888260}{}
\TC{tea-club-25.mp3}{\TCtitle{\#25}{Dimension Brothers}}{2016-11-20}{\Luke, \MBerlinger}{Jakou hru mám hrát s přítelkyní nebo rodiči? Taková otázka není až tak vzácná. Položil si jí i Michal Berlinger, a protože je programátor, sám si odpověděl hrou Dimension Brothers. V lokální kooperativní skákačce si dva hráči pomáhají překonat neviditelné překážky, které ale musí nejprve sami objevit. Kvůli červeným a modrým brýlím totiž vidí každý z nich herní svět jinak. Michal v Tea Clubu nepovídá jen o prototypování, zastaví se i u problematiky přístupnosti lokálních multiplayerů a zabrousí na půdu gamejamů.}{0:21:51}{10107644}{}
\TC{tea-club-26.mp3}{\TCtitle{\#26}{Vikings – Wolves of Midgard}}{2016-11-26}{\Luke, \PNagy}{Když vám vypálí vesnici a máte štěstí, tak utečete. Když máte smůlu, tak ne. A když mají smůlu oni, tak vypálili vesnici klanu Ulfungů a je úplně jedno, že jsou to obři a mají velké kyje a temnou magii. Takový je námět akční RPG od slovenského studia Games Farm, jehož šéfa Petera Nagyho jsme zastihli s kamerou a mikrofonem během pražské Game Developers Session. Peter nám toho o běsnících vikinzích vyprávěl skoro celou ságu!}{0:18:17}{10601168}{}
\TC{tea-club-27.mp3}{\TCtitle{\#27}{Children of the Galaxy}}{2017-04-14}{\Luke, \FDusek}{Již čtyři roky vyvíjí Filip Dušek 4X strategii Children of the Galaxy, ve které se obrací na velikány žánru jako je Stallaris, Galactic Civilization nebo Endless Space. Zároveň přidává řadu systémů, které objevování, bitvy a kolonizování osvěžují. Jeden příklad za všechny: Tahové souboje s jednotkami, kterým se v průběhu tažení vylepšují statistiky. Pro Filipa jde sice o první vlastní hru, ale zkušenosti nasbíral již ve společnosti 2K a studiem žánrových spřízněnců. V našem Tea Clubu se zmiňuje i o tom, proč je výhodné dělat programátorské video streamy.}{0:30:31}{15628278}{}
\TC{tea-club-28.mp3}{\TCtitle{\#28}{LostHero}}{2018-07-09}{\Luke, }{Povídání o vývoji Lost Hero}{0:24:55}{10768712}{}
% \TC{tea-club-29.mp3}{\TCtitle{\#29}{}}{2018--}{\Luke, }{}{}{}{}